\section*{Capítulo \ref{ch:g11:vect} --- Vectores y rectas en el plano}

\begin{solution}{exo:g11:vect:1}
$\vect{AB}\,(3, 4)$, $\vect{AC}\,(-3, 2)$,
$\vect{AB} + \vect{AC}\,(0, 6)$,
$2\vect{AB} - 3\vect{AC}\,(6 + 9,\ 8 - 6) = (15, 2)$. El \omterm{prop:g10:coordgeom:midpoint}{punto medio} de
$[BC]$ es $\left(\frac{5 + (-1)}{2}, \frac{3+1}{2}\right) = (2, 2)$.
\end{solution}

\begin{solution}{exo:g11:vect:2}
$\det = 4 \times 3 - (-2) \times (-6) = 12 - 12 = 0$: \omterm{def:g11:vect:collinear}{colineales} (en
efecto, $\vec v = -\frac12 \vec u$).

$\det = 2 \times 7 - 3 \times 5 = -1 \neq 0$: no \omterm{def:g11:vect:collinear}{colineales}.

$\det = 0 \times (-8) - 0 \times 3 = 0$: \omterm{def:g11:vect:collinear}{colineales} (los dos son
verticales).
\end{solution}

\begin{solution}{exo:g11:vect:3}
Por $A(2,1)$ con dirección $(3, -1)$: identificando $(-b, a) = (3, -1)$ se
obtiene $a = -1$, $b = -3$, luego $-x - 3y + c = 0$; sustituyendo $A$:
$-2 - 3 + c = 0$, $c = 5$. Multiplicando por $-1$: $x + 3y - 5 = 0$.

Por $B(0,4)$ y $C(2,0)$: $\vect{BC}\,(2, -4)$, luego $a = -4$ y $b = -2$:
$-4x - 2y + c = 0$; sustituyendo $B$: $-8 + c = 0$, $c = 8$.
Simplificando entre $-2$: $2x + y - 4 = 0$. Comprobación con $C$:
$4 + 0 - 4 = 0$.
\end{solution}

\begin{solution}{exo:g11:vect:4}
Un \omterm{def:g11:vect:direction}{vector director} es $(-b, a) = (2, 3)$. Cortes: $x = 0$ da $y = 3$,
punto $(0, 3)$; $y = 0$ da $x = -2$, punto $(-2, 0)$. Para $P(2,6)$:
$3 \times 2 - 2 \times 6 + 6 = 0$, luego $P \in d$. Para $Q(1,4)$:
$3 - 8 + 6 = 1 \neq 0$, luego $Q \notin d$.
\end{solution}

\begin{solution}{exo:g11:vect:5}
$\vect{AB}\,(3, 6)$ y $\vect{AC}\,(-3, -6)$:
$\det = 3 \times (-6) - (-3) \times 6 = 0$, luego $A$, $B$, $C$ están
alineados.

$\vect{DE}\,(2, 3)$ y $\vect{DF}\,(5, 8)$:
$\det = 2 \times 8 - 5 \times 3 = 1 \neq 0$: $D$, $E$, $F$ no están
alineados.
\end{solution}

\begin{solution}{exo:g11:vect:6}
Primera pareja: sumando las dos ecuaciones,
$(2x + y - 7) + (x - y + 1) = 3x - 6 = 0$, luego $x = 2$ y después
$y = x + 1 = 3$: se cortan en $(2, 3)$ (comprobación: $4 + 3 - 7 = 0$).

Segunda pareja: $ab' - a'b = 1 \times 4 - (-2) \times (-2) = 0$, así que
las rectas son paralelas; multiplicando $e$ por $-2$ queda
$-2x + 4y - 6 = 0$, distinta de $e'$ ($-6 \neq 1$): estrictamente
paralelas, sin punto de corte.
\end{solution}

\begin{solution}{exo:g11:vect:7}
Para $d: 2x - y + 3 = 0$: dirección $(-b, a) = (1, 2)$, y el punto
$(0, 3)$ está en $d$, así que $d = \{(t,\ 3 + 2t) : t \in \R\}$.

Para la recta paramétrica que pasa por $(1, -3)$ con dirección $(2, 1)$:
$a = 1$, $b = -2$, luego $x - 2y + c = 0$; sustituyendo $(1, -3)$:
$1 + 6 + c = 0$, $c = -7$. \omterm{def:g10:algebra:equation}{Ecuación}: $x - 2y - 7 = 0$ (comprobación con
$t = 1$, punto $(3, -2)$: $3 + 4 - 7 = 0$).
\end{solution}

\begin{solution}{exo:g11:vect:8}
Una recta paralela tiene los mismos $(a, b)$: $4x - 3y + c = 0$; pasando
por $P(1, -2)$: $4 + 6 + c = 0$, luego $c = -10$ y la recta es
$4x - 3y - 10 = 0$.

Que $mx + 2y - 5 = 0$ sea paralela a $d$ exige
$ab' - a'b = 4 \times 2 - m \times (-3) = 8 + 3m = 0$, luego
$m = -\frac83$.
\end{solution}

\begin{solution}{exo:g11:vect:9}
El \omterm{prop:g10:coordgeom:midpoint}{punto medio} de $[AB]$ es
$M\!\left(\frac{-1+3}{2}, \frac{0+2}{2}\right) = (1, 1)$. La \omterm{def:g10:stats:median}{mediana} es la
recta $(CM)$ con $C(1, -4)$: los dos puntos tienen \omterm{def:g10:coordgeom:system}{abscisa} $1$, así que la
\omterm{def:g10:stats:median}{mediana} es la recta vertical
\[
x - 1 = 0 .
\]
\end{solution}

\begin{solution}{exo:g11:vect:10}
$\det(\vec u, \vec v) = m(m - 1) - 3 \times 2 = m^2 - m - 6$. Se anula
cuando $m^2 - m - 6 = 0$: $\Delta = 1 + 24 = 25$, raíces
$m = \frac{1 \pm 5}{2}$, es decir, $m = 3$ o $m = -2$.
\end{solution}

\begin{solution}{exo:g11:vect:11}
\emph{1.} $I = \left(\frac12, 0\right)$. Entonces
$\vect{DI} = \left(\frac12 - 0,\ 0 - 1\right) = \left(\frac12, -1\right)$
y $\vect{DJ} = \frac23\vect{DI} = \left(\frac13, -\frac23\right)$, luego
\[
J = D + \vect{DJ} = \left(\frac13,\ \frac13\right).
\]

\emph{2.} $\vect{AJ}\left(\frac13, \frac13\right)$ y $\vect{AC}\,(1, 1)$:
$\det = \frac13 \times 1 - 1 \times \frac13 = 0$, así que $A$, $J$, $C$
están alineados; de hecho, $\vect{AJ} = \frac13\vect{AC}$: la recta $(DI)$
corta la diagonal $(AC)$ exactamente a un tercio de su longitud desde
$A$.
\end{solution}

\begin{solution}{pb:g11:vect:1}
\textbf{1.} $x = 1 + 3t$, $y = 2 - t$. De la segunda, $t = 2 - y$;
sustituyendo: $x = 1 + 3(2 - y) = 7 - 3y$: \omterm{def:g10:algebra:equation}{ecuación} general
$x + 3y - 7 = 0$.

\textbf{2.} Para $ax + by + c = 0$, un \omterm{def:g11:vect:direction}{vector director} es $(-b, a)$: aquí
$(5, 2)$. Un punto: $y = 1$ da $x = 1$, o sea, $(1, 1)$. Paramétrica:
$x = 1 + 5t$, $y = 1 + 2t$.

\textbf{3.} Sustituyendo $x = 7 - 3y$ en $2x - 5y + 3 = 0$:
$14 - 6y - 5y + 3 = 0$, luego $y = \frac{17}{11}$ y
$x = 7 - \frac{51}{11} = \frac{26}{11}$: el punto
$\left(\frac{26}{11}, \frac{17}{11}\right)$.

\textbf{4.} Direcciones $(5, 2)$ y $(3, -1)$:
$\det = 5 \times (-1) - 2 \times 3 = -11 \neq 0$: se cortan, como se ha
visto. Para la segunda pareja: las direcciones $(5, 2)$ y $(10, 4)$ tienen
$\det = 0$ (paralelas) y, doblando $2x - 5y + 3 = 0$, queda
$-4x + 10y - 6 = 0 \neq -4x + 10y + 1 = 0$: estrictamente paralelas.

\textbf{5.} $\vect{AB}(3, 2)$, $\vect{AC}(9, 6)$:
$\det = 18 - 18 = 0$: alineados.

\textbf{6.} $P$: de $x = 2t$, $y = 3 + t$ se obtiene
$y = 3 + \frac x2$. $Q$: de $x = 10 - t$, $t = 10 - x$, luego
$y = 2(10 - x) - 1 = 19 - 2x$. Cruce: $3 + \frac x2 = 19 - 2x$ da
$x = 6.4$, $y = 6.2$: las trayectorias se cruzan en $(6.4,\ 6.2)$.

\textbf{7.} $P$ pasa por allí cuando $2t = 6.4$: $t = 3.2$ h. $Q$, cuando
$10 - t = 6.4$: $t = 3.6$ h. Veinticuatro minutos de \omterm{def:g11:seq:arithmetic}{diferencia}: no hay
colisión. Advertencia: una carta náutica muestra \emph{trayectorias}; solo
los relojes paramétricos dicen si dos movimientos ocupan el mismo punto en
el mismo instante.

\textbf{8.} $D^2(t) = (3t - 10)^2 + (4 - t)^2 = 10t^2 - 68t + 116$:
vértice en $t = \frac{68}{20} = 3.4$ h, donde $D^2 = 0.4$: la máxima
\omterm{def:g10:numbers:approx}{aproximación} es $D = \sqrt{0.4} \approx 0.63$ millas en $t = 3.4$.

\textbf{9.} Se incumple: $0.63 < 1$. $D^2(t) < 1$ se lee
$10t^2 - 68t + 115 < 0$: $\Delta = 4624 - 4600 = 24$, raíces
$\frac{68 \pm \sqrt{24}}{20} = 3.4 \pm 0.245$: los barcos están demasiado
cerca desde $t \approx 3.16$ hasta $t \approx 3.64$, unos $29$ minutos de
infracción. Uno de los dos tiene que cambiar de rumbo.

\textbf{10.} Cada coordenada de cada barco es una \omterm{def:g11:func:function}{función} \emph{afín} de
$t$, así que cada \omterm{def:g11:seq:arithmetic}{diferencia} de \omterm{def:g10:coordgeom:system}{coordenadas} es afín y $D^2$, suma de dos
cuadrados de funciones afines, es de segundo grado (con coeficiente
principal no negativo). Por eso la máxima \omterm{def:g10:numbers:approx}{aproximación} se lee siempre en
un vértice.

\textbf{11.} Encontrarse con $P$ en $t = 3$, es decir, en $(6, 6)$:
$(3a, 3b) = (6, 6)$ da $(a, b) = (2, 2)$, con velocidad
$\sqrt{8} \approx 2.8$ nudos.

\textbf{12.} $D(0, 1)$, $I\left(\frac12, 0\right)$: dirección
$\left(\frac12, -1\right)$, es decir, $(1, -2)$: \omterm{def:g10:algebra:equation}{ecuación} $2x + y = 1$.

\textbf{13.} $(AC)$: $y = x$. Entonces $2x + x = 1$: $x = \frac13$: el
punto $\left(\frac13, \frac13\right)$, a un tercio del camino de $A$ a
$C$.

\textbf{14.} $B(1, 0)$, $K\left(\frac12, 1\right)$: dirección
$\left(-\frac12, 1\right)$, \omterm{def:g10:algebra:equation}{ecuación} $2x + y = 2$. Con $y = x$:
$x = \frac23$: el punto $\left(\frac23, \frac23\right)$. Las dos cevianas
cortan $(AC)$ en sus dos puntos de trisección: teorema demostrado, dos
veces tres líneas de aritmética.

\textbf{15.} Cualquier paralelogramo se lleva sobre el cuadrado unidad
tomando $A$ como origen y $\vect{AB}$ y $\vect{AD}$ como unidades de los
dos ejes. Esa elección conserva todo aquello de lo que habla el teorema
(\omterm{prop:g10:coordgeom:midpoint}{puntos medios}, alineación y razones sobre una misma recta, todo
expresable con \omterm{def:g11:vect:vector}{vectores} y escalares), así que demostrar el enunciado en el
cuadrado lo demuestra para todo paralelogramo. (En cambio, \emph{no}
transportaría longitudes ni ángulos: eso sí lo deforma el sistema.)

\textbf{16.} $E\left(1, \frac12\right)$: la recta $(AE)$ es
$y = \frac x2$. La diagonal $(DB)$, de $(0,1)$ a $(1,0)$, es
$y = 1 - x$. Corte: $\frac x2 = 1 - x$: $x = \frac23$: el punto
$\left(\frac23, \frac13\right)$, otra vez un punto de trisección, ahora de
la otra diagonal. A las cevianas que salen de \omterm{prop:g10:coordgeom:midpoint}{puntos medios} les encantan
los tercios.

\textbf{17.} Las dos primeras: $x + y = 4$ y $2x - y = 2$ dan $3x = 6$:
$(2, 2)$. La tercera: $2 - 4 = -2$: se cumple. Son concurrentes en
$(2, 2)$.

\textbf{18.} $mx - y + 2 - 3m = m(x - 3) + (2 - y) = 0$: para que la
igualdad valga para \emph{todo} $m$, hay que tomar $x = 3$ e $y = 2$; y,
en efecto, todas las $d_m$ contienen a $(3, 2)$. Es un haz de rectas por
un punto, una recta por cada \omterm{def:g10:reffunc:affine}{pendiente} $m$.

\textbf{19.} Paralela a $y = 2x + 7$: \omterm{def:g10:reffunc:affine}{pendiente} $m = 2$, la recta
$2x - y - 4 = 0$. Por el origen: $2 - 3m = 0$: $m = \frac23$.

\textbf{20.} General: sustituyes un punto y la pertenencia queda
resuelta. Explícita: \omterm{def:g10:reffunc:affine}{pendiente} y \omterm{def:g10:reffunc:affine}{ordenada en el origen} a la vista, \omterm{def:g10:functions:graph}{gráfica}
inmediata. Paramétrica: el reloj va incorporado, con colisiones, máxima
\omterm{def:g10:numbers:approx}{aproximación} e intercepción. El \omterm{def:g11:vect:det}{determinante}: un producto cruzado y el
paralelismo queda decidido. El sistema de referencia bien elegido: un
paralelogramo se vuelve cuadrado unidad y un teorema clásico de trisección
se vuelve tres cortes de rectas.
\end{solution}
